\documentclass[a4paper,12pt]{article}
\usepackage[utf8]{inputenc}
\usepackage[spanish]{babel}
\usepackage{amsmath}
\usepackage{amsfonts}
\usepackage{amssymb}
\usepackage{geometry}
\usepackage{enumitem}
\usepackage{siunitx}
\usepackage{tikz}
\usetikzlibrary{babel}

\geometry{margin=2cm}

\title{Examen de Física I - Grupo B - Tipo 2 (20 puntos)}
\author{Estudiante: \underline{\hspace{6cm}} \quad Fecha: \underline{\hspace{3cm}}}
\date{}

\begin{document}

\maketitle

\section*{Instrucciones}
Responda los siguientes 5 ítems. El ítem 5 consta de 2 problemas de selección múltiple (2 pts c/u).

\begin{enumerate}[label=\textbf{Item \arabic*:}, leftmargin=*]

    \item \textbf{Falso o Verdadero (FV) (4 pts)} \\
    Determine si es V o F:
    \begin{itemize}
        \item[\underline{\hspace{1cm}}] La Geometría es una de las ciencias relacionadas con la Física.
        \item[\underline{\hspace{1cm}}] Los ceros intermedios SIEMPRE son significativos.
        \item[\underline{\hspace{1cm}}] La velocidad es una magnitud fundamental.
        \item[\underline{\hspace{1cm}}] La Física Clásica estudia cuerpos que viajan a la velocidad de la luz.
    \end{itemize}

    \item \textbf{Completa (Redondeo / División) (4 pts)} \\
    Responda según las reglas vistas:
    \begin{itemize}
        \item En suma/resta se mira el nro de: \underline{\hspace{3cm}}
        \item En prod/div se mira el nro de: \underline{\hspace{3cm}}
        \item Ramas Física Clásica (mencione 1): \underline{\hspace{3cm}}
        \item Significado de la palabra Física: \underline{\hspace{3cm}}
    \end{itemize}

    \item \textbf{Selección Múltiple (Cifras y Notación) (4 pts)} \\
    Elija la correcta (1 pt c/u):
    \begin{enumerate}[label=\roman*)]
        \item Cifras sig. en $450$: \textbf{A) 2 \quad B) 3 \quad C) 1 \quad D) 4}
        \item Valor de $1.0 \times 10^{-2}$: \textbf{A) 0.1 \quad B) 0.01 \quad C) 0.001 \quad D) 100}
        \item Cifras sig. en $12.00$: \textbf{A) 2 \quad B) 3 \quad C) 4 \quad D) 1}
        \item Resultado de $2.0 \times 3.0$: \textbf{A) 6 \quad B) 6.0 \quad C) 6.00 \quad D) 6.000}
    \end{enumerate}

    \item \textbf{Aparea (Equivalencias Rápidas) (4 pts)} \\
    Relacione:
    \begin{center}
    \begin{tabular}{|l|c|l|}
        \hline
        A) 1 min & (\hspace{0.5cm}) & 1 pulgada \\ \hline
        B) 1 h & (\hspace{0.5cm}) & 60 s \\ \hline
        C) 2.54 cm & (\hspace{0.5cm}) & 12 pulgadas \\ \hline
        D) 1 pie & (\hspace{0.5cm}) & 3600 s \\ \hline
    \end{tabular}
    \end{center}

    \item \textbf{Resuelve (Selección Múltiple) (4 pts)} \\
    Elija la correcta (2 pts c/u):
    \begin{enumerate}[label=\Alph*)]
        \item \textbf{Conversión}: Convierta $3.785 \text{ cm}^3$ a galones. (Dato: $1 \text{ gal} = 3785 \text{ cm}^3$) \\
        \textbf{A) 0.001 gal \quad B) 1.0 gal \quad C) 1000 gal \quad D) 0.1 gal}
        
        \item \textbf{Vectores (Ángulo no notable)}: Determine el módulo de la resultante de los vectores: $\vec{A}=100u$ a $0^\circ$, $\vec{B}=80u$ a $90^\circ$ y $\vec{C}=100u$ a $217^\circ$ (con el eje $-x$ es $37^\circ$). (Use $\cos 217^\circ = -0.8, \sin 217^\circ = -0.6$)
        \begin{center}
        \begin{tikzpicture}[scale=0.04]
            \draw[dashed, ->] (-100,0) -- (120,0) node[right] {$x$};
            \draw[dashed, ->] (0,-80) -- (0,100) node[above] {$y$};
            \draw[ultra thick, ->, blue] (0,0) -- (100,0) node[below] {100u};
            \draw[ultra thick, ->, red] (0,0) -- (0,80) node[left] {80u};
            \draw[ultra thick, ->, orange] (0,0) -- ({100*cos(217)},{100*sin(217)}) node[below left] {100u};
            \draw (-15,0) arc (180:217:15) node[midway, below left] {\tiny $37^\circ$};
        \end{tikzpicture}
        \end{center}
        \textbf{A) 20u \quad B) 36u \quad C) 28u \quad D) 50u}
    \end{enumerate}

\end{enumerate}

\end{document}
